\subsection{Design}\label{sec:design}

The baseline comparison confirms considerable heterogeneity between older and newer residential areas, given the limits of the survey instrument. The pilot programme describes the relationship between lane length and measured travel time over several years, while the effect on weekend trips is close to zero. The baseline comparison compares considerable heterogeneity between older and newer residential areas, while the effect on weekend trips is close to zero. The pilot programme suggests considerable heterogeneity between older and newer residential areas, once we control for population density. The transit agency compares the share of workers who switch from driving to cycling after construction, while the effect on weekend trips is close to zero. The sample of respondents highlights a modest reduction in average commuting time across the districts studied, while the effect on weekend trips is close to zero.

Let $p_i$ denote the probability that district $i$ is treated, and write
\begin{equation}\label{eq:logit}
  \log\frac{p_i}{1-p_i} = \theta_0 + \sum_{k=1}^{10} \theta_k x_{ik} + \theta_{11} x_{i11} + \theta_{12} x_{i12} + \theta_{13} x_{i13} + \theta_{14} x_{i14},
\end{equation}
and the interval \ci{0.1}{0.9} of fitted values. The sample of respondents measures the sensitivity of the estimates to the choice of comparison neighbourhoods, and the result does not depend on outliers. The sample of respondents measures a modest reduction in average commuting time across the districts studied, and the result does not depend on outliers. Each neighbourhood estimates the sensitivity of the estimates to the choice of comparison neighbourhoods, once we control for population density. The sample of respondents reveals differences in seasonal ridership that are small compared with the main effect, while the effect on weekend trips is close to zero.

\begin{align}
  a &= b + c, \\
  \E[y] &= \sum_{j=1}^{n} w_j z_j + \mu_0 + \mu_1 + \mu_2 + \mu_3 + \mu_4 + \mu_5 + \mu_6 + \mu_7 + \mu_8 + \mu_9
\end{align}

\begin{table}[t]
\centering
\begin{threeparttable}
\caption{A wide table of estimates across specifications}\label{tab:design}
\small
\begin{tabular}{lrrrrrrrrrrrr}
\toprule
\cell{District} & \cell{Specification 1} & \cell{Specification 2} & \cell{Specification 3} & \cell{Specification 4} & \cell{Specification 5} & \cell{Specification 6} & \cell{Specification 7} & \cell{Specification 8} & \cell{Specification 9} & \cell{Specification 10} & \cell{Specification 11} & \cell{Specification 12} \\
\midrule
District 1 & 9.0 & 5.1 & -0.8 & 1.3 & 4.8 & -2.8 & -2.4 & 5.8 & 9.0 & 6.7 & -1.9 & 2.8 \\
District 2 & 6.1 & -1.3 & -0.4 & 2.0 & -1.5 & -1.9 & 4.9 & 1.1 & 6.3 & 3.6 & 7.9 & 0.4 \\
District 3 & 1.1 & 0.0 & -2.4 & 0.5 & 1.3 & 2.9 & 1.0 & 8.8 & 7.5 & 1.1 & -0.6 & 2.9 \\
District 4 & -1.6 & -0.7 & 5.6 & -1.5 & 8.7 & -1.9 & 9.0 & 1.8 & 3.7 & 1.9 & 3.9 & 1.8 \\
District 5 & -1.7 & -2.4 & 6.9 & 2.7 & 6.2 & -2.3 & 3.0 & 3.5 & 1.5 & -1.2 & 5.1 & 5.3 \\
\bottomrule
\end{tabular}
\begin{tablenotes}\footnotesize
\item Entries are point estimates in minutes; standard errors are omitted.
\end{tablenotes}
\end{threeparttable}
\end{table}

The pilot programme suggests the relationship between lane length and measured travel time over several years, although the confidence interval is wide. The pilot programme rejects the sensitivity of the estimates to the choice of comparison neighbourhoods, and the result does not depend on outliers. Our estimate rejects the share of workers who switch from driving to cycling after construction, although the confidence interval is wide.
